\section{Three finite algebraic systems}\label{sec:systems}

Theorem~\ref{thm:complete} says that, in the covered cases, a
nonconstant meromorphic solution is one of three kinds and has one
pole modulo its periods. Each kind is a finite-dimensional family of
functions: a polynomial in \(1/z\), a polynomial in a fixed simply
periodic function with one pole, or a constant plus a linear
combination of \(\wp\) and its derivatives. Substituting the general member of each family
into \eqref{eq:main} therefore gives finitely many polynomial
equations in its coefficients. This section writes down the three
systems, fixes their normalizations, and proves that they are exactly
equivalent to the one-pole solutions (Theorem~\ref{thm:systems}).
They are instances of the principal-part construction of
\citet{DeminaKudryashov2011,DeminaKudryashovElliptic2011}.

The systems all have the same shape. Let \(V\) be a polynomial
expression in some variables, and let \(\delta\) be a differentiation
rule for those variables, so that the polynomials in them form a
differential algebra with derivation \(\delta\). Write
\[
 \mathcal E_P(V;\delta)=
 \delta^pV+\sum_{j=0}^{p-1}a_j\delta^jV+\frac12V^2.
\]
Each system below is the polynomial identity
\(\mathcal E_P(V;\delta)=0\) for an appropriate \(V\) and \(\delta\).
A polynomial identity means the vanishing of every coefficient.

\subsection{Rational solutions}

A rational solution whose only pole is at zero, of order \(p\), and
which has no nonconstant polynomial part at infinity is a polynomial
of degree \(p\) in \(T=1/z\). Since \(DT=-T^2\), the equation becomes
\begin{equation}\label{eq:rational-system}
 V_R(T)=b+\sum_{k=1}^{p}r_kT^k,\qquad r_p=c_\ast,\qquad
 \mathcal E_P(V_R;-T^2\partial_T)=0.
\end{equation}
The left side of the identity is a polynomial of degree at most
\(2p\) in \(T\), so this is a finite system in
\(b,r_1,\ldots,r_{p-1}\).
No rational solution has a nonconstant polynomial part at infinity:
such a part would have its degree doubled by \(v^2\), while the
linear differential expression cannot reach that degree.
Consequently \eqref{eq:rational-system} is the full rational
one-pole system.

\subsection{Rational functions of an exponential}

For \(q=\kappa^2\ne0\), define
\begin{equation}\label{eq:T}
 T_q(z)=\frac{\kappa}{2}\coth\frac{\kappa z}{2},
 \qquad
 DT_q=\frac q4-T_q^2.
\end{equation}
The expression is independent of the choice of sign of \(\kappa\).
It has period \(2\pi i/\kappa\) and one pole per period, at zero,
with residue one; it plays the role that \(1/z\) plays in the
rational case.
The finite system is
\begin{equation}\label{eq:exponential-system}
 V_S(T)=\sum_{k=0}^{p}u_kT^k,\qquad
 u_p=c_\ast,\qquad q\ne0,\qquad
 \mathcal E_P\left(V_S;\left(\frac q4-T^2\right)\partial_T\right)=0.
\end{equation}

To see that every one-pole solution rational in an exponential has
this form, write \(v=R(t)\), \(t=e^{\kappa z}\).
The rational function \(R\) has finite limits at both \(t=0\) and
\(t=\infty\). If it had a pole at either end, an extreme Laurent
monomial would be doubled by \(R^2\) and could not cancel against
\(P(\kappa t\partial_t)R\).
After a translation, its only finite pole is at \(t=1\), of order \(p\).
Such a function is a degree-\(p\) polynomial in \(T_q\).
The logistic coordinate
\[
 Q=\frac1{1+e^{-z}},\qquad DQ=Q(1-Q),
\]
is an equivalent coordinate after fixing the rate and translating the
pole; \(Q(z+i\pi)=T_1(z)+1/2\).
Polynomial logistic representations and their use in constructing
travelling waves are developed by
\citet{Kudryashov2012,Kudryashov2013,Kudryashov2015}.

The rational system is the \(q=0\) specialization of
\eqref{eq:exponential-system}, with \(T_0=1/z\).
We keep \(q=0\) and \(q\ne0\) distinct so that this degeneration is
never counted twice.

\subsection{Elliptic solutions}

Use the Weierstrass algebra
\begin{equation}\label{eq:weierstrass}
 Y^2=4X^3-g_2X-g_3,\qquad
 \mathscr D X=Y,\qquad
 \mathscr D Y=6X^2-\frac{g_2}{2},\qquad
 \Delta=g_2^3-27g_3^2\ne0.
\end{equation}
Its analytic realization is \(X=\wp(z;g_2,g_3)\),
\(Y=\wp'(z;g_2,g_3)\); see \citet[Chapter~23]{DLMF}.
Set
\begin{equation}\label{eq:elliptic-system}
 V_E=b+\sum_{j=0}^{p-2}t_j\mathscr D^jX,\qquad
 t_{p-2}=-s_p,\qquad
 s_p=\frac{2(2p-1)!}{((p-1)!)^2},
 \qquad
 \mathcal E_P(V_E;\mathscr D)=0.
\end{equation}
Using \eqref{eq:weierstrass}, every polynomial in \(X\) and \(Y\)
reduces to the normal form \(F(X)+YG(X)\). The system consists of
the coefficients of \(F\) and \(G\) set equal to zero; all of them
are polynomials with rational coefficients in the displayed unknowns.

The representation is the classical one. An elliptic function with
one pole has zero residue there.
The functions \(1,\wp,\wp',\ldots,\wp^{(p-2)}\) have distinct
principal parts and span the space of elliptic functions whose only
pole is at zero, of order at most \(p\).
Subtracting the indicated principal part from a one-pole solution
leaves an entire elliptic function, hence a constant.
The leading coefficient in \eqref{eq:elliptic-system} follows from
\[
 D^{p-2}\wp=(-1)^{p-2}(p-1)!z^{-p}+O(1).
\]
The nonsingularity condition \(\Delta\ne0\) is retained throughout;
the singular cubics correspond to the rational and simply periodic
systems, which we have already written down.

\begin{theorem}[Finite matching and exhaustion]\label{thm:systems}
For every \(p\ge2\), the systems
\eqref{eq:rational-system}, \eqref{eq:exponential-system}, and
\eqref{eq:elliptic-system} are necessary and sufficient for the
respective one-pole solutions of \eqref{eq:main}, with the pole at
zero and with the stated normalizations.
For the operator classes of Theorem~\ref{thm:complete}, these systems,
together with the constant solutions, exhaust all globally meromorphic
solutions.

The systems describe exceptional parameter values and continuous
families as well as isolated solutions. Their sufficiency follows
from substitution, without any genericity assumption.
\end{theorem}

\begin{proof}
Necessity is the principal-part argument given for each kind above.
Conversely, a solution of a coefficient system gives a globally
meromorphic function by its displayed realization, and the chain rule
gives \(E_P(v)=0\).
The leading coefficient \(c_\ast\ne0\) ensures a nonconstant function.
In the elliptic case every nonsingular pair \((g_2,g_3)\) has a
Weierstrass realization, and the normal form \(F(X)+YG(X)\) is unique.
The exhaustion statement is Theorem~\ref{thm:complete}.
\end{proof}

\subsection{Eliminating the profile coefficients}

The systems above treat the operator and the profile coefficients as
unknowns together, which is the form the enumeration in
Section~\ref{sec:counting} requires. For a fixed operator they can be
reduced further, because the Laurent recurrence already determines
the principal part of any solution. Compute
\[
 r_k=c_{p-k}\quad(1\le k\le p),\qquad s=c_p
\]
from \eqref{eq:recurrence}.
These coefficients occur before the even-order resonance and are
determined by \(P\) at either parity.
In the simply periodic algebra define
\[
 S_k(T;q)=\frac{(-1)^{k-1}}{(k-1)!}
 \left[\left(\frac q4-T^2\right)\partial_T\right]^{k-1}T.
\]
It has principal part \(z^{-k}\).
Its regular Laurent constant is
\[
 d_k(q)=
 \begin{cases}
 0,&k\text{ odd},\\
 -B_kq^{k/2}/k!,&k\text{ even},
 \end{cases}
\]
where \(B_k\) denotes a Bernoulli number.
Thus the unique candidate with the required principal part and
constant is
\begin{equation}\label{eq:genuszero-candidate}
 V_0(T;P,q)=s+\sum_{k=1}^{p}r_k\bigl(S_k(T;q)-d_k(q)\bigr).
\end{equation}
Substituting it into \eqref{eq:main} leaves polynomial equations in
the single unknown \(q\).

For the elliptic candidate, first require \(r_1=0\), since an
elliptic function with one pole has no residue there, and put
\begin{equation}\label{eq:elliptic-candidate}
 \widetilde V_E=
 s+\sum_{k=2}^{p}r_k
 \left(\frac{(-1)^{k-2}}{(k-1)!}\mathscr D^{\,k-2}X
       -e_k(g_2,g_3)\right),
\end{equation}
where \(e_k\) is the regular Laurent constant of the corresponding
derivative of \(\wp\).
It is a polynomial in \(g_2,g_3\), computed from
\eqref{eq:weierstrass}.
The remaining equations involve only \(g_2,g_3\), with
\(\Delta\ne0\).
This is the coefficient-comparison construction used in
\citet{DeminaKudryashov2011,DeminaKudryashov2014}.

For even \(P\) the equation \(E_P(v)=0\) does not see the energy: a
candidate may solve it on any energy level. To classify a fixed
level, compute \(h\) from \eqref{eq:energyh} for the prescribed
energy and append
\begin{equation}\label{eq:selecth}
 [z^{2p}]V_0(T_q;P,q)=h
 \quad\text{or}\quad
 [z^{2p}]\widetilde V_E(\wp,\wp';P,g_2,g_3)=h.
\end{equation}
Both are polynomial equations obtained from finite Laurent expansions.
With this equation included, the reduced systems classify a fixed
energy level: at fixed \(P\) of odd degree, or fixed \((P,h)\) in the
even case, there is at most one nonconstant solution with its pole at
zero.

When all prescribed data are exact algebraic numbers, that is, the
coefficients of \(P\) and, in the even case, the energy, these are
finitely many polynomial equations in one or two unknowns.
Polynomial gcds in the one-variable case and elimination in the
multivariable case decide whether a solution exists.
A nonvanishing condition such as \(\Delta\ne0\) is imposed by
adjoining \(u\Delta-1=0\). Elimination with case distinctions
describes instead the operators, or the marked data \((P,h)\), that
admit each kind of solution, as a finite union of systems of
polynomial equations and nonvanishing conditions
\citep{CoxLittleOShea2005}.
